\documentclass[11pt, reqno]{article}

\usepackage{amsmath, amsthm, amssymb}
\usepackage{enumitem}
\usepackage{tcolorbox}
\usepackage{hyperref}
\usepackage{tikz}
\usepackage{tikz-cd}
\usepackage{pgfplots}
\pgfplotsset{compat=1.18}
\usetikzlibrary{arrows.meta}
\usepackage{mathrsfs}
\usepackage{fancyhdr}
\usepackage[bottom=0.75in, top=1in, left=0.5in, right=0.5in]{geometry}
\usepackage{array}   % for \newcolumntype macro
\newcolumntype{L}{>{$}l<{$}}

\theoremstyle{plain}
\newtheorem*{theorem}{Theorem}
\newtheorem*{proposition}{Proposition}
\newtheorem{exercise}{Exercise}
\newtheorem*{lemma}{Lemma}
\newtheorem*{corollary}{Corollary}

\theoremstyle{definition}
\newtheorem*{definition}{Definition}
\newtheorem*{example}{Example}

\theoremstyle{remark}
\newtheorem*{remark}{Remark}

\renewcommand{\phi}{\varphi}
\renewcommand{\epsilon}{\varepsilon}
\renewcommand{\emptyset}{\varnothing}

\newcommand{\RR}{\mathbb{R}}
\newcommand{\ZZ}{\mathbb{Z}}
\newcommand{\NN}{\mathbb{N}}
\newcommand{\CC}{\mathbb{C}}
\newcommand{\QQ}{\mathbb{Q}}

\DeclareMathOperator{\ima}{\text{im}}

\begin{document}

\topmargin=-40pt
\rhead{Henry Woodburn}
\lhead{Math 622}
\renewcommand{\headrulewidth}{1pt}
\renewcommand{\headsep}{20pt}
\thispagestyle{fancy}

{\Huge \bfseries \noindent Homework 6}

\begin{enumerate}
    \item[1.] We show that 
    \[
        V_1^* \otimes \cdots V_k^* \otimes W \simeq \text{Mult}(V_1\times \cdots \times V_k, W),
    \]
    where $\text{Mult}(V_1\times \cdots \times V_k, W)$ is the space of multilinear maps $V_1 \times \cdots \times V_k
    \to W$.

    We proved in class the isomorphism 
    \[
        V_1^*\otimes \cdots\otimes V_k^* \simeq (V_1 \otimes \cdots \otimes V_k)^*,
    \]
    so that
    \[
        V_1^*\otimes \cdots \otimes V_k^*\otimes W \simeq 
        (V_1^*\otimes \cdots \otimes V_k^*)\otimes W \simeq (V_1 \otimes \cdots \otimes V_k)^*\otimes W.
    \]

    Moreover, we proved in class that $U^*\otimes V \simeq \text{Hom}(U, V)$. Applying this above, we have 
    \[
        (V_1 \otimes \cdots \otimes V_k)^*\otimes W \simeq \text{Hom}(V_1 \otimes\cdots\otimes V_k, W).
    \]
    To finish, we must show $\text{Hom}(V_1 \otimes\cdots\otimes V_k, W) \simeq \text{Mult}(V_1\times \cdots \times V_k, W)$.

    For each $\phi \in \text{Mult}(V_1\times \cdots \times V_k, W)$, by the universal property of the tensor
    product, there is a unique map $\tilde{\phi} \in \text{Hom}(V_1 \otimes\cdots\otimes V_k, W)$ such
    that $\tilde{\phi}(v_1 \otimes \cdots \otimes v_k) = \phi(v_1, \dots, v_k)$. Let 
    $\Phi$ be the map $\phi \mapsto \tilde{\phi}$. It is injective by the uniqueness mentioned above. 
    
    To show surjectivity, let $g \in \text{Hom}(V_1 \otimes\cdots\otimes V_k, W)$. Then 
    the map $h(v_1,\dots, v_k) = g(v_1\otimes \cdots \otimes v_k)$ is multilinear and satisfies $\Phi(h) = g$. 

    $\Phi$ is also a linear map, since for $g, h \in \text{Mult}(V_1\times \cdots \times V_k, W)$,
    \begin{align*}
        \Phi(g + h)(v_1\otimes \cdots \otimes v_k) &= (g + h)(v_1, v_2,\dots, v_k)
        = \alpha g(v_1, \dots, v_k) + h(v_1, \dots, v_k) \\
        &= \Phi(g)(v_1\otimes \cdots\otimes v_k) + \Phi(h)(v_1\otimes \cdots\otimes v_k).
    \end{align*}
    Then we have shown the isomorphism.

    \item[2.] 
    
\end{enumerate}

\end{document}